\section{Methods and experimental protocol}

\subsection{Models and datasets}

\headingindent
The experiments used ResNet-18, ResNet-50, and ViT-B/16 teachers. ResNet
students consumed flattened feature maps, globally averaged features, or
selected feature coordinates. ViT feature-denoising students used CLS
embeddings or patch-token grids. The student family included linear models,
multilayer perceptrons, random forests, and residual convolutional denoisers.

The first evaluation used CIFAR-10 and CIFAR-100. The
opposite CIFAR dataset was Near-OOD, while MNIST and SVHN were Far-OOD. The
last phase adopted the OpenOOD v1.5 benchmark \cite{zhang2023openood}. OpenOOD
is an effort to standardize OOD evaluation through fixed datasets, image-list
splits, and aggregation groups. Many papers about OOD detection often use
different OOD datasets or different splits under the same benchmark name,
making their reported metrics impossible to compare directly. OpenOOD v1.5
added Tiny ImageNet (TIN), Textures, and Places365 to the CIFAR protocol, and
introduced ImageNet-200 with SSB-hard and NINCO as Near-OOD and iNaturalist,
Textures, and OpenImage-O as Far-OOD.

\begin{table}[ht]
\centering
\rowcolors{2}{rowblue!35}{white}
\begin{tabularx}{\textwidth}{@{}L{2.2cm}L{2.9cm}YY@{}}
\toprule
Phase & ID dataset & Near-OOD & Far-OOD \\
\midrule
Initial CIFAR & CIFAR-10 & CIFAR-100 & MNIST, SVHN \\
Initial CIFAR & CIFAR-100 & CIFAR-10 & MNIST, SVHN \\
OpenOOD CIFAR-10 & CIFAR-10 & CIFAR-100, TIN & MNIST, SVHN, Textures,
Places365 \\
OpenOOD CIFAR-100 & CIFAR-100 & CIFAR-10, TIN & MNIST, SVHN, Textures,
Places365 \\
OpenOOD ImageNet & ImageNet-200 & SSB-hard, NINCO & iNaturalist, Textures,
OpenImage-O \\
\bottomrule
\end{tabularx}
\caption{Dataset protocols used during the internship.}
\label{tab:dataset-protocols}
\end{table}

\subsection{Baseline distillation}

\headingindent
The baseline student receives a clean teacher feature and predicts teacher
logits or probabilities. We compared linear, MLP, and random-forest students.
This establishes whether limited capacity alone creates an OOD signal before
any explicit corruption is added.

\subsection{Embedding perturbations}

\headingindent
Embedding perturbations modify an intermediate teacher activation before the
student sees it. The tested operations were percentile clipping, Monte Carlo
dropout, PCA projection, and masked PCA projection. Clipping was applied to a
single internal teacher layer or sequentially to several internal teacher layers. Its threshold could be global,
channel-dependent, or spatially dependent. Training targets were either the
clean teacher output or the teacher output after the same perturbation.
The student received perturbed feature vectors concatenated with perturbation parameters.

\subsection{Pixel perturbations}

\headingindent
Pixel perturbation was tested with affine and photometric augmentation and
with PixMix. The affine pipeline sampled rotation, translation, scaling,
brightness, and contrast. The student received a feature extracted from the
perturbed image concatenated with normalized transformation
parameters. PixMix repeatedly combined a crop-and-flip view with an
augmentation or an external structurally complex image. These experiments
tested whether prediction stability under input changes separates ID and OOD
examples.

\subsection{Subspace-student ensemble}

\headingindent
The ensemble trained 16 linear students on channel or whitened PCA subspaces.
One variant used fixed ordered partitions; another used seeded random
partitions. In both cases, the 512 feature coordinates were divided among the
students so that each coordinate was assigned to exactly one student.
Predictive entropy and BALD measured uncertainty and disagreement across the
16 predictions \cite{houlsby2011bald}.

% \subsection{Selection and metrics}

% \headingindent
% Early reports selected a configuration by
% \[
%   C=\tfrac12\left(\mathrm{macro\ AUROC}
%   +1-\mathrm{macro\ FPR95}\right).
% \]
% Later reports first averaged within the Near-OOD and Far-OOD groups, then gave
% each group half of the macro weight. Results in the next sections are reported
% as \textbf{AUROC / FPR@95}. Higher AUROC and lower FPR@95 are better. Section
% \ref{sec:metric-conventions} distinguishes the two FPR@95 definitions that
% were encountered near the end of the internship.
